\begin{document}
\section{Figures}
\begin{figure}[t]
  \centering
  \includegraphics[width=0.8\linewidth]{plots/loss}
  \caption{Training loss over time.}
  \label{fig:loss}
\end{figure}

\begin{figure*}[!ht]
  \centering
  \includegraphics[height=3cm, angle=90]{images/arch.pdf}
  \caption*{Architecture without a number.}
\end{figure*}

\begin{figure}
  \centering
  \begin{subfigure}{0.45\textwidth}
    \includegraphics[width=\linewidth]{a/left}
    \caption{Left panel.}
    \label{fig:left}
  \end{subfigure}
  \hfill
  \begin{subfigure}{0.45\textwidth}
    \includegraphics{a/right.png}
    \caption{Right panel.}
    \label{fig:right}
  \end{subfigure}
  \caption{Two panels.}
  \label{fig:pair}
\end{figure}

\begin{figure}
  \begin{tikzpicture}[scale=2]
    \draw (0,0) -- (1,1);
  \end{tikzpicture}
  \caption{A TikZ diagram.}
\end{figure}

Figure~\ref{fig:loss} and \ref{fig:pair} show results.
An inline image \includegraphics{icon} stays inline.
\end{document}
